\SetPageTheme{story}
\PageHeader{Lectia 3}{Ce se intampla la puteri ale lui 10}{L3 • P1}

\begin{missionbox}[Mai mult decat o cifra]
  Daca impartirea la 10 ne ajuta sa lucram cu ultima cifra, atunci puterile lui 10 ne ajuta sa lucram cu ultimele doua, trei sau mai multe cifre.

  Aceasta idee este utila nu doar in probleme scolare. In prelucrarea informatiei, uneori trebuie sa izolam capete de date, grupuri de simboluri sau bucati fixe dintr-un cod.
\end{missionbox}

\begin{conceptbox}[Generalizarea]
  \begin{itemize}
    \item \texttt{n \% 100} da ultimele doua cifre;
    \item \texttt{n / 100} elimina ultimele doua cifre;
    \item \texttt{n \% 1000} da ultimele trei cifre;
    \item \texttt{n / 1000} elimina ultimele trei cifre.
  \end{itemize}
\end{conceptbox}

\begin{guidebox}[Inceput ghidat]
  Pentru \texttt{n = 5384}, ultimele doua cifre sunt \hrulefill iar numarul ramas dupa \texttt{/ 100} este \hrulefill
\end{guidebox}

\newpage
\SetPageTheme{guided}
\PageHeader{Lectia 3}{Vedem mecanismul}{L3 • P2}

\begin{examplebox}[Exemple explicate]
  \begin{itemize}
    \item \texttt{5384 \% 100 = 84}
    \item \texttt{5384 / 100 = 53}
    \item \texttt{5384 \% 1000 = 384}
    \item \texttt{5384 / 1000 = 5}
  \end{itemize}
\end{examplebox}

\begin{guidebox}[Exercitii ghidate]
  Completeaza:
  \begin{itemize}
    \item \texttt{76429 \% 100 =} \hrulefill
    \item \texttt{76429 / 100 =} \hrulefill
    \item \texttt{76429 \% 1000 =} \hrulefill
    \item \texttt{76429 / 1000 =} \hrulefill
  \end{itemize}
\end{guidebox}

\newpage
\SetPageTheme{guided}
\PageHeader{Lectia 3}{Comparatii utile}{L3 • P3}

\begin{solobox}[Exercitii semi-ghidate]
  1. Completeaza tabelul:
  \begin{center}
  \begin{tabularx}{\linewidth}{|X|X|X|X|X|}
    \hline
    \textbf{n} & \textbf{\% 10} & \textbf{/ 10} & \textbf{\% 100} & \textbf{/ 100} \\
    \hline
    9062 & & & & \\
    \hline
    430 & & & & \\
    \hline
  \end{tabularx}
  \end{center}

  2. Scrie in cuvinte diferenta dintre \texttt{\% 10} si \texttt{\% 100}.
\end{solobox}

\begin{solobox}[Independent scurt]
  Da un exemplu de situatie in care ai vrea sa izolezi ultimele trei cifre ale unui numar.
  \AnswerSpace{3}
\end{solobox}

\newpage
\SetPageTheme{challenge}
\PageHeader{Lectia 3}{Consolidare independenta}{L3 • P4}

\begin{solobox}[Fisa independenta]
  1. Calculeaza pentru \texttt{n = 902718}:
  \begin{itemize}
    \item \texttt{n \% 10}
    \item \texttt{n / 10}
    \item \texttt{n \% 100}
    \item \texttt{n / 100}
    \item \texttt{n \% 1000}
    \item \texttt{n / 1000}
  \end{itemize}

  2. Explica de ce puterile lui 10 sunt potrivite pentru lucrul cu cifrele.
  \AnswerSpace{8}
\end{solobox}
